\scanpage{80}
toute la structure étant décrite canoniquement en termes du bitorseur $\operatorname{Hom}_{\widehat\Pi_1(X)}(x,y)$ sous $(\widehat\Pi_1(X,y),\widehat\Pi_1(X,x))$ et l'opération de $\Gamma$ sur \deleted{tout} ce fourbis. On se rappellera surtout que $f_{\unsure{y,x}}$ et $g_{y,x}$ se déterminent mutuellement, et qu'ils sont déterminés par la connaissance d'\uline{un seul} \deleted{$f_l$ \inserted{$:\Gamma\mathbin{-}\widehat\Pi_1(X,x)$} \unsure{(ou}} \deleted{$g_l$ (i.e. $f_l(u)$)} $f_l$, ou $g_l$
\[
\tag{33}\left\{\begin{aligned}
f_l&:\Gamma\unsure{\mathbin{-}|\longrightarrow}\widehat\Pi_1(X,x)\\
g_l&:\Gamma\unsure{\mathbin{-}|\longrightarrow}\widehat\Pi_1(\unsure{Y|X},y)
\end{aligned}\right.
\qquad l:x\mathbin{-}y\quad(\text{dans }\Pi_1(X))
\]
(Si on a \deleted{\unsure{un}} un $T^\alpha$ $1$-connexe $\subset X$, contenant $x,y$) ils sont connus quand on connaît seulement $g^{T^\alpha}_{x,y}$.

\margin{\uline{NB} $f_{x,y}$ et $f_{y,x}$, ou $g_{x,y}$ et $g_{y,x}$, se déterminent mutuellement, car, en termes du chemin $l$, on a
\[
\tag{34}\left\{\begin{aligned}
g_{l^{-1}}(u)&=f_l(u)^{-1}\\
f_{l^{-1}}(u)&=g_l(u)^{-1}
\end{aligned}\right.
\]}

Quand $X$ se plonge dans une VA, plus généralement s'il se plonge dans une extension d'une VA par un tore (p.ex. $X$ courbe alg. ``\unsure{anabélienne|abélienne}''), alors le th. de Mordell-Weil implique que pour un $x_0$ \inserted{$\in X_K$} fixé (comme pt de référence --- mais ce $x_0$ \deleted{\unsure{vient}} joue un rôle bidon!)
\margin{79}
